\chapter{类型1休假到达过程的概率生成函数PGF推导}
\label{app:proof_lemma1}

本附录推导类型1休假期内新到达数据包量的概率生成函数PGF。在类型1休假期的第一个时隙，可能出现三种情形。

\begin{itemize}
	\item \textbf{情形1}：以概率$qe^{-\theta}$，标记节点向接收机传输数据包，其他$n-1$个节点均不请求传输。此时休假期只持续一个时隙。一个时隙内到达量服从参数为$\lambda$的伯努利分布，其概率生成函数为$I(z)=1-\lambda+\lambda z$。因此，$A_{V1}$在该情形下的概率生成函数为
	\begin{equation}
		G_{A_{V1}}(z)|\text{C1} = 1 - \lambda + \lambda z
	\end{equation}
	
	\item \textbf{情形2}：以概率$(1-q)\theta e^{-\theta}$，标记节点不传输，其他$n-1$个节点中有且仅有一个成功传输。标记节点继续处于休假期。此后，它将依次经历三个阶段：(1) 一个非竞争的时隙；(2) 另一个节点的忙期；(3) 一个新的类型1休假期。因此，情况2下$A_{V1}$的概率生成函数为：
	\begin{equation}
		G_{A_{V1}}(z)|\text{C2} = (1 - \lambda + \lambda z) \times G_{A_{B}}(z) \times G_{A_{V1}}(z)
	\end{equation}
	
	\item \textbf{情形3}：以概率$1-(1-q)\theta e^{-\theta}-qe^{-\theta}$，第一个时隙中无节点成功。标记节点在下一个时隙继续竞争信道。此后，它将经历两个阶段：(1) 一个空闲时隙；(2) 一个类型1的休假期。因此，$A_{V1}$在该情形下的概率生成函数为
	\begin{equation}
		G_{A_{V1}}(z)|\text{C3} = (1 - \lambda + \lambda z) \times G_{A_{V1}}(z)
	\end{equation}
\end{itemize}

基于全概率定律，综合以上三种情形，我们即可得到目标结果。
